\documentclass[11pt]{article}
\usepackage[margin=1in]{geometry}
\usepackage{amsmath,amssymb}
\usepackage{enumitem}
\usepackage[T1]{fontenc}

\title{COMP2300 Set 5 Practice Exam Solutions}
\author{}
\date{}

\begin{document}
\maketitle

\section*{Part A: Multiple Choice}
\begin{enumerate}[label=\textbf{A\arabic*.}]
  \item B
  \item C
  \item B
  \item D
  \item C
  \item B
  \item B
  \item A
  \item B
  \item B
\end{enumerate}

\section*{Part B: Computational Problems}
\begin{enumerate}[label=\textbf{B\arabic*.}]
  \item \textbf{State sequences and cycle counts}
  \begin{enumerate}[label=(\alph*)]
    \item
    \begin{itemize}[leftmargin=1.5em]
      \item \texttt{LDR}: \texttt{Fetch -> Decode -> MemAdr -> MemRead -> MemWB}, so 5 cycles.
      \item \texttt{STR}: \texttt{Fetch -> Decode -> MemAdr -> MemWrite}, so 4 cycles.
      \item DP immediate: \texttt{Fetch -> Decode -> ExecuteI -> ALUWB}, so 4 cycles.
      \item \texttt{B}: \texttt{Fetch -> Decode -> Branch}, so 3 cycles.
    \end{itemize}
    \item The 6-instruction sequence takes:
    \[
      5 + 4 + 4 + 3 + 4 + 5 = 25 \text{ cycles}
    \]
    Average CPI:
    \[
      \frac{25}{6} \approx 4.17
    \]
  \end{enumerate}

  \item \textbf{PC, R15, and branch target calculation}
  \begin{enumerate}[label=(\alph*)]
    \item Reading \texttt{R15} should supply \texttt{PC + 8}, so:
    \[
      0x80A0 + 8 = 0x80A8
    \]
    \item
    \[
      \texttt{ExtImm} = \text{sign-extend}(0x000003 \ll 2) = 0x0000000C
    \]
    \[
      \texttt{BTA} = \texttt{PC+8} + \texttt{ExtImm} = 0x80A8 + 0x0000000C = 0x80B4
    \]
  \end{enumerate}

  \item \textbf{Decoder outputs}
  \begin{center}
  \begin{tabular}{c|c|c|c}
  Instruction class & \texttt{Op} & \texttt{RegSrc1:0} & \texttt{ImmSrc} \\ \hline
  \texttt{LDR} & 01 & 10 & 01 \\
  \texttt{STR} & 01 & 10 & 01 \\
  DP immediate & 00 & 00 & 00 \\
  DP register & 00 & 00 & 00 \\
  \texttt{B} & 10 & 01 & 10 \\
  \end{tabular}
  \end{center}
  Explanation:
  \begin{itemize}[leftmargin=1.5em]
    \item \texttt{RegSrc0 = 1} only when \texttt{Op = 10}.
    \item \texttt{RegSrc1 = 1} only when \texttt{Op = 01}.
    \item \texttt{ImmSrc = Op}.
  \end{itemize}

  \item \textbf{Average CPI and total cycles}
  \begin{enumerate}[label=(\alph*)]
    \item
    \[
      \text{Average CPI} = (0.25)(5) + (0.10)(4) + (0.13)(3) + (0.52)(4)
    \]
    \[
      = 1.25 + 0.40 + 0.39 + 2.08 = 4.12
    \]
    \item Total cycles for 50 million instructions:
    \[
      50 \times 10^6 \times 4.12 = 206 \times 10^6
    \]
    So the total is 206 million cycles.
  \end{enumerate}

  \item \textbf{Critical path and clock frequency}
  \begin{enumerate}[label=(\alph*)]
    \item First compute:
    \[
      t_{ALU} + t_{mux} = 120 + 25 = 145 \text{ ps}
    \]
    Hence:
    \[
      \max(145, 200) = 200 \text{ ps}
    \]
    Therefore:
    \[
      T_{c2} = 40 + 2(25) + 200 + 50 = 340 \text{ ps}
    \]
    \item
    \[
      f_{max} = \frac{1}{340 \times 10^{-12}} \approx 2.94 \times 10^9 \text{ Hz}
    \]
    So:
    \[
      f_{max} \approx 2.94 \text{ GHz}
    \]
  \end{enumerate}

  \item \textbf{Execution time}
  \begin{enumerate}[label=(\alph*)]
    \item
    \[
      T = N \times CPI \times T_c
    \]
    \[
      = (100 \times 10^9)(4.12)(340 \times 10^{-12})
    \]
    \[
      = 140.08 \text{ s} \approx 140 \text{ s}
    \]
    \item Revised design:
    \[
      T' = (100 \times 10^9)(4.12)(300 \times 10^{-12}) = 123.6 \text{ s}
    \]
    \item Speedup:
    \[
      \text{Speedup} = \frac{140.08}{123.6} \approx 1.13
    \]
  \end{enumerate}
\end{enumerate}

\section*{Part C: Past Exam Questions}
\begin{enumerate}[label=\textbf{C\arabic*.}]
  \item \textbf{ISA vs Microarchitecture}
  \begin{enumerate}[label=(\alph*)]
    \item Ripple-carry vs carry-save adder: \textbf{Microarchitecture}
    \item Instruction width: \textbf{ISA}
    \item Branch prediction: \textbf{Microarchitecture}
    \item Hardwired vs microprogrammed control: \textbf{Microarchitecture}
  \end{enumerate}

  \item \textbf{Multi-Cycle Execution Time}
  \\[0.5ex]
  Average CPI:
  \[
    (0.50)(6) + (0.25)(8) + (0.25)(4) = 3 + 2 + 1 = 6
  \]
  Execution time:
  \[
    T = \frac{2 \times 10^9 \times 6}{1 \times 10^9} = 12 \text{ s}
  \]

  \item \textbf{Single-Cycle vs Pipelining}
  \\[0.5ex]
  Two drawbacks of a single-cycle processor are:
  \begin{itemize}[leftmargin=1.5em]
    \item The clock period must be long enough for the slowest instruction.
    \item Hardware must be provisioned for worst-case resource use, so utilization is poor.
  \end{itemize}
  Pipelining addresses these by overlapping different stages of different instructions, increasing throughput and making better use of otherwise idle hardware.\\
  ILP is the ability to execute or make progress on multiple independent instructions at the same time. Two microarchitecture types that exploit ILP are \textbf{superscalar} and \textbf{VLIW}.

  \item \textbf{ILP Ranking}
  \\[0.5ex]
  Highest to lowest ILP potential:
  \[
    \text{A} > \text{B} > \text{C}
  \]
  where:
  \begin{itemize}[leftmargin=1.5em]
    \item A = scalar in-order pipeline with forwarding and hazard detection
    \item B = scalar in-order pipeline without forwarding or hazard detection
    \item C = multi-cycle processor
  \end{itemize}
  Justification:
  \begin{itemize}[leftmargin=1.5em]
    \item A can overlap instructions while also reducing stalls via forwarding and hazard handling.
    \item B still overlaps instructions, but extra NOPs and stalls reduce usable ILP.
    \item C has the least ILP because it does not pipeline instruction stages across multiple instructions.
  \end{itemize}
\end{enumerate}

\end{document}
